\section{Conclusion}
\label{sec:conclusion}

We presented \textsc{Graphify}, an open-source system for multi-modal knowledge
graph construction from heterogeneous software corpora. Its hybrid two-pass
extraction pipeline combines deterministic tree-sitter AST parsing across 33
language variants with LLM semantic extraction across 6 file modalities, tagging
every edge with a confidence class for verifiability. Embedding-free Leiden
community detection recovers subsystem structure without any vector model.
Hub-aware BFS/DFS traversal delivers a median of 68.6$\times$ fewer tokens per
query than naive full-corpus reading, with per-question range of
43.5$\times$--126.7$\times$ on the 52-file multi-modal corpus. The case study in
Section~\ref{sec:case} demonstrates that the graph surfaces cross-modal and
cross-repository connections --- such as code nodes linked to the research papers
describing their algorithms --- that static analysis cannot produce.

Graphify is deployed across 21 AI assistant platforms and has accumulated 1.2M+
PyPI downloads, providing real-world evidence that pre-built knowledge graphs
are a practical replacement for reactive file-reading in LLM coding assistants.

\smallskip
\noindent\textbf{Future work} includes: (1) cross-repository graph federation
under a shared schema; (2) extending the heterogeneous-corpus KG model to
continuous personal knowledge --- meetings, emails, and browser history pose
novel challenges around temporal knowledge evolution and privacy-preserving
graph construction; (3) replacing the $p_{99}$ hub-suppression threshold with a
degree centrality classifier trained on community cohesion signals.

\section*{Acknowledgments}
The author thanks the open-source community for contributions, bug reports,
and real-world testing across 20+ countries that shaped the system's design.
The author discloses affiliation with Y Combinator (S26 cohort).
